\documentclass[conference]{IEEEtran}

% ============================================================
% PAQUETES
% ============================================================

\usepackage[utf8]{inputenc}
\usepackage[T1]{fontenc}
\usepackage[spanish]{babel}

\usepackage{amsmath}
\usepackage{amssymb}

\usepackage{graphicx}
\usepackage{float}
\usepackage{booktabs}
\usepackage{array}

\usepackage{cite}
\usepackage{url}

\usepackage{xcolor}


% ============================================================
% INFORMACIÓN DEL DOCUMENTO
% ============================================================

\title{
Análisis espectral de señales de audio mediante la Transformada Rápida de Fourier para la identificación de sonidos animales, instrumentos musicales y voces humanas
}

\author{
\IEEEauthorblockN{
Yuber Garnica,
Miguel Gordillo,
Nombre Apellido
}

\IEEEauthorblockA{
Programa de Ingeniería Mecatrónica\\
Universidad Militar Nueva Granada\\
Bogotá, Colombia\\
correo@ejemplo.com
}
}


% ============================================================
% DOCUMENTO
% ============================================================

\begin{document}

\maketitle


% ============================================================
% RESUMEN
% ============================================================

\begin{abstract}

En este trabajo se realiza un análisis espectral de diferentes señales de audio mediante la Transformada Rápida de Fourier (FFT). Se estudiaron tres categorías de señales: sonidos producidos por animales, sonidos correspondientes a instrumentos musicales y voces humanas. Para cada señal se obtuvo su representación en el dominio temporal y frecuencial, además de calcular la media, la desviación estándar y la frecuencia dominante. Los resultados permitieron observar diferencias en la distribución espectral de las señales y determinar qué características presentan mayor utilidad para su diferenciación. Se encontró que la frecuencia dominante y la distribución general del espectro proporcionan información más útil para distinguir las señales que la media y la desviación estándar. Sin embargo, debido al reducido número de muestras utilizadas, los resultados no permiten establecer conclusiones estadísticas generalizables.

\end{abstract}


% ============================================================
% PALABRAS CLAVE
% ============================================================

\begin{IEEEkeywords}

Transformada Rápida de Fourier, FFT, procesamiento digital de señales, análisis espectral, señales de audio, frecuencia dominante.

\end{IEEEkeywords}


% ============================================================
% I. INTRODUCCIÓN
% ============================================================

\section{Introducción}

Las señales de audio constituyen un tipo de señal temporal cuya información puede analizarse tanto en el dominio del tiempo como en el dominio de la frecuencia. Mientras que la representación temporal permite observar cómo varía la amplitud de una señal a lo largo del tiempo, el análisis frecuencial permite identificar los componentes espectrales que la conforman.

Una de las herramientas fundamentales para realizar este análisis es la Transformada de Fourier, la cual permite representar una señal en términos de sus componentes frecuenciales. En aplicaciones digitales, el cálculo directo de la Transformada Discreta de Fourier puede resultar computacionalmente costoso, por lo que se utiliza la Transformada Rápida de Fourier (FFT), un algoritmo eficiente para obtener la representación espectral de una señal.

En el presente trabajo se analiza un conjunto de señales de audio correspondientes a tres animales, tres instrumentos musicales y dos personas. Para cada señal se obtuvo su representación temporal y su espectro mediante FFT. Adicionalmente, se calcularon parámetros estadísticos como la media y la desviación estándar.

El objetivo principal es determinar si las características obtenidas permiten diferenciar las señales pertenecientes a cada categoría y establecer cuáles de estas características presentan mayor utilidad para su identificación.


% ============================================================
% II. FUNDAMENTO TEÓRICO
% ============================================================

\section{Fundamento teórico}


\subsection{Señales de audio}

Una señal de audio puede representarse como una función de amplitud dependiente del tiempo, denotada como $x(t)$. En sistemas digitales, una señal analógica es muestreada a una frecuencia determinada para obtener una secuencia discreta:

\begin{equation}
x[n] = x(nT_s)
\end{equation}

donde $T_s$ representa el periodo de muestreo.

En este trabajo se utilizaron señales con una frecuencia de muestreo de $48\,000$ Hz. De acuerdo con el teorema de Nyquist, la máxima frecuencia representable sin aliasing es aproximadamente:

\begin{equation}
f_{\max} = \frac{f_s}{2}
\end{equation}

Por lo tanto:

\begin{equation}
f_{\max} = \frac{48000}{2} = 24000\text{ Hz}
\end{equation}


\subsection{Transformada de Fourier}

La Transformada de Fourier permite representar una señal en función de sus componentes frecuenciales. Para una señal continua puede expresarse como:

\begin{equation}
X(f) =
\int_{-\infty}^{\infty}
x(t)e^{-j2\pi ft}dt
\end{equation}

Esta transformación permite analizar qué frecuencias están presentes en una señal y con qué magnitud.


\subsection{Transformada Discreta de Fourier}

Para señales digitales se utiliza la Transformada Discreta de Fourier (DFT):

\begin{equation}
X[k] =
\sum_{n=0}^{N-1}
x[n]e^{-j2\pi kn/N}
\end{equation}

donde $N$ representa el número de muestras.


\subsection{Transformada Rápida de Fourier}

La Transformada Rápida de Fourier (FFT) es un algoritmo utilizado para calcular eficientemente la DFT. En el presente trabajo se empleó la función \texttt{rfft()} de la biblioteca NumPy de Python.

La magnitud del espectro se obtuvo mediante:

\begin{equation}
|X[k]| =
\sqrt{
\operatorname{Re}(X[k])^2+
\operatorname{Im}(X[k])^2
}
\end{equation}

La frecuencia dominante se definió como aquella frecuencia asociada al máximo valor de magnitud del espectro.


\subsection{Media y desviación estándar}

La media de una señal discreta se calculó mediante:

\begin{equation}
\mu =
\frac{1}{N}
\sum_{n=1}^{N}x[n]
\end{equation}

La desviación estándar se calculó mediante:

\begin{equation}
\sigma =
\sqrt{
\frac{1}{N}
\sum_{n=1}^{N}
(x[n]-\mu)^2
}
\end{equation}

Estas medidas permiten caracterizar estadísticamente la amplitud de las señales.


% ============================================================
% III. METODOLOGÍA
% ============================================================

\section{Metodología}

El procedimiento experimental se dividió en tres grupos: animales, instrumentos musicales y voces humanas.

Cada archivo de audio fue almacenado en formato WAV y procesado mediante Python utilizando las bibliotecas NumPy, SciPy y Matplotlib.

El procedimiento general fue:

\begin{enumerate}

\item Cargar el archivo de audio.

\item Identificar la frecuencia de muestreo y el número de muestras.

\item Convertir las señales estéreo a mono cuando fue necesario.

\item Limitar la duración de las señales a un máximo de 10 segundos.

\item Representar la señal en el dominio temporal.

\item Calcular la FFT.

\item Obtener el espectro de magnitud.

\item Determinar la frecuencia dominante.

\item Calcular la media y la desviación estándar.

\item Comparar los resultados obtenidos entre las diferentes señales.

\end{enumerate}


\subsection{Señales animales}

Se analizaron tres sonidos correspondientes a un gato, un perro y un gallo.

Los archivos utilizados fueron:

\begin{itemize}

\item Gato.wav
\item Perro.wav
\item Gallo.wav

\end{itemize}


% ============================================================
% FIGURAS ANIMALES
% ============================================================

\begin{figure}[H]
\centering
\includegraphics[width=\linewidth]{imagen_2026-08-18_160028246.png}
\caption{Señal temporal y espectro frecuencial correspondiente al sonido de un gato.}
\label{fig:gato}
\end{figure}


\begin{figure}[H]
\centering
\includegraphics[width=\linewidth]{imagen_2026-08-18_160957543.png}
\caption{Señal temporal y espectro frecuencial correspondiente al sonido de un perro.}
\label{fig:perro}
\end{figure}


\begin{figure}[H]
\centering
\includegraphics[width=\linewidth]{imagen_2026-08-18_161847927.png}
\caption{Señal temporal y espectro frecuencial correspondiente al sonido de un gallo.}
\label{fig:gallo}
\end{figure}


\begin{figure}[H]
\centering
\includegraphics[width=\linewidth]{imagen_2026-08-18_161952064.png}
\caption{Comparación de los espectros frecuenciales de los sonidos animales.}
\label{fig:animales_comparacion}
\end{figure}


% ============================================================
% INSTRUMENTOS
% ============================================================

\subsection{Instrumentos musicales}

Para estudiar diferentes registros frecuenciales se seleccionaron tres instrumentos musicales:

\begin{itemize}

\item Contrabajo, como instrumento de baja frecuencia.

\item Piano, como instrumento de frecuencia media.

\item Flauta, como instrumento de frecuencia alta.

\end{itemize}


\begin{figure}[H]
\centering
\includegraphics[width=\linewidth]{imagen_2026-08-19_080009805.png}
\caption{Señal temporal y espectro frecuencial correspondiente al contrabajo.}
\label{fig:contrabajo}
\end{figure}


\begin{figure}[H]
\centering
\includegraphics[width=\linewidth]{imagen_2026-08-19_080152644.png}
\caption{Señal temporal y espectro frecuencial correspondiente al piano.}
\label{fig:piano}
\end{figure}


\begin{figure}[H]
\centering
\includegraphics[width=\linewidth]{imagen_2026-08-19_080253090.png}
\caption{Señal temporal y espectro frecuencial correspondiente a la flauta.}
\label{fig:flauta}
\end{figure}


\begin{figure}[H]
\centering
\includegraphics[width=\linewidth]{imagen_2026-08-19_080432552.png}
\caption{Comparación de los espectros frecuenciales de los instrumentos musicales.}
\label{fig:instrumentos_comparacion}
\end{figure}


% ============================================================
% VOCES
% ============================================================

\subsection{Voces humanas}

Finalmente, se analizaron dos grabaciones correspondientes a dos personas diferentes. Ambas personas pronunciaron la misma frase con el propósito de controlar el contenido lingüístico de las señales.

La frase utilizada fue:

\begin{quote}
``Hoy estamos analizando señales de voz mediante la transformada de Fourier.''
\end{quote}


\begin{figure}[H]
\centering
\includegraphics[width=\linewidth]{imagen_2026-08-19_080432552.png}
\caption{Señal temporal y espectro frecuencial correspondiente a la Persona 1.}
\label{fig:persona1}
\end{figure}


\begin{figure}[H]
\centering
\includegraphics[width=\linewidth]{persona2.png}
\caption{Señal temporal y espectro frecuencial correspondiente a la Persona 2.}
\label{fig:persona2}
\end{figure}


\begin{figure}[H]
\centering
\includegraphics[width=\linewidth]{personas_comparacion.png}
\caption{Comparación de los espectros frecuenciales de las dos voces humanas.}
\label{fig:personas_comparacion}
\end{figure}


% ============================================================
% IV. RESULTADOS
% ============================================================

\section{Resultados}

Los resultados obtenidos mediante el procesamiento de las señales se resumen en la Tabla~\ref{tab:resultados}.


\begin{table}[!t]
\centering
\caption{Resultados del análisis de las señales}
\label{tab:resultados}

\small

\begin{tabular}{@{}l l r r r@{}}
\toprule
Categoría & Señal & Media & Desv. Est. & $f_d$ (Hz) \\
\midrule

Animal & Gato & 0.036 & 2288.47 & 1940.90 \\
Animal & Perro & -0.107 & 2047.43 & 491.00 \\
Animal & Gallo & 0.049 & 3006.66 & 1334.90 \\

\midrule

Instrumento & Contrabajo & -1.244 & 3836.46 & 146.48 \\
Instrumento & Piano & -0.101 & 1745.00 & 384.69 \\
Instrumento & Flauta & 0.585 & 3826.29 & 432.00 \\

\midrule

Voz & Persona 1 & 6.699 & 2778.54 & 214.30 \\
Voz & Persona 2 & -0.167 & 2153.47 & 142.47 \\

\bottomrule
\end{tabular}

\end{table}


\subsection{Resultados de los sonidos animales}

En los sonidos animales se obtuvieron frecuencias dominantes de $1940.90$ Hz para el gato, $491.00$ Hz para el perro y $1334.90$ Hz para el gallo.

La diferencia entre las frecuencias dominantes permitió observar una separación espectral entre las tres muestras. El perro presentó la frecuencia dominante más baja, mientras que el gato presentó la más alta.

La comparación de los espectros puede observarse en la Fig.~\ref{fig:animales_comparacion}.


\subsection{Resultados de los instrumentos}

Para los instrumentos se obtuvieron frecuencias dominantes de $146.48$ Hz para el contrabajo, $384.69$ Hz para el piano y $432.00$ Hz para la flauta.

El contrabajo presentó claramente la frecuencia dominante más baja. El piano y la flauta presentaron valores superiores, aunque relativamente cercanos entre sí.

Es importante considerar que los instrumentos musicales no producen una única frecuencia, sino una frecuencia fundamental acompañada de componentes armónicos. Por esta razón, el espectro completo proporciona información adicional a la frecuencia dominante.


\subsection{Resultados de las voces}

Las frecuencias dominantes obtenidas para las dos personas fueron de $214.30$ Hz y $142.47$ Hz, respectivamente.

Aunque ambas personas pronunciaron la misma frase, se observaron diferencias en sus espectros frecuenciales. Esto demuestra que las características de la voz generan diferencias observables en el dominio de la frecuencia.

Sin embargo, una sola grabación por persona no es suficiente para establecer que estas características permitan identificar de manera general a un individuo.


% ============================================================
% V. DISCUSIÓN
% ============================================================

\section{Discusión}

Los resultados muestran que la Transformada Rápida de Fourier constituye una herramienta útil para analizar señales de audio y observar diferencias entre diferentes fuentes sonoras.

La frecuencia dominante fue la característica que presentó mayor utilidad para distinguir las muestras analizadas. En el caso de los animales, las frecuencias dominantes estuvieron suficientemente separadas para observar diferencias claras entre las tres señales.

En el caso de los instrumentos, el contrabajo presentó una frecuencia dominante considerablemente menor que el piano y la flauta. Sin embargo, la diferencia entre piano y flauta fue menor, lo que demuestra que la clasificación de un instrumento como grave, medio o agudo depende de la nota o muestra específica utilizada.

Por otra parte, la media de las señales presentó valores cercanos a cero en la mayoría de los casos. Esto indica que dicha característica no resulta adecuada para diferenciar las fuentes sonoras.

La desviación estándar presentó diferencias entre las señales, pero no mostró una relación directa con la frecuencia dominante. Por ejemplo, el contrabajo y la flauta presentaron desviaciones estándar similares a pesar de tener características espectrales diferentes. Por esta razón, la desviación estándar tampoco puede utilizarse como único criterio de identificación.

Debe considerarse además que la magnitud de la FFT depende de factores como la amplitud de grabación, la duración de la señal y las condiciones de captura. Por lo tanto, la magnitud máxima no debe utilizarse directamente como criterio de identificación sin realizar previamente una normalización.

Finalmente, aunque se observaron diferencias entre las señales, el número reducido de muestras impide determinar si dichas diferencias son estadísticamente significativas o generalizables. Para desarrollar un sistema de identificación más robusto sería necesario utilizar múltiples grabaciones de cada categoría y considerar características adicionales como la frecuencia fundamental, distribución de armónicos, energía por bandas y otras características espectrales.


% ============================================================
% VI. CONCLUSIONES
% ============================================================

\section{Conclusiones}

La Transformada Rápida de Fourier permitió transformar las señales de audio del dominio temporal al dominio frecuencial, facilitando la identificación de componentes espectrales característicos de cada muestra.

En los sonidos animales se observaron diferencias claras entre las frecuencias dominantes del gato, perro y gallo. Para las muestras utilizadas, estas diferencias permitieron distinguir las señales mediante su comportamiento espectral.

En los instrumentos musicales se verificó una tendencia creciente en la frecuencia dominante desde el contrabajo hasta la flauta. Sin embargo, la separación entre el piano y la flauta fue menor que la existente entre el contrabajo y los otros instrumentos.

En las voces humanas también se observaron diferencias espectrales entre las dos personas, a pesar de que ambas pronunciaron la misma frase. Esto demuestra que las características individuales de la voz se reflejan en el dominio frecuencial.

La media y la desviación estándar no demostraron ser características suficientes para identificar por sí solas las diferentes señales. En contraste, la frecuencia dominante y el análisis del espectro completo proporcionaron información más relevante para la diferenciación.

Finalmente, los resultados obtenidos deben considerarse exploratorios debido al reducido número de muestras. Para determinar si las diferencias encontradas son estadísticamente significativas sería necesario utilizar un conjunto de datos mayor, con varias grabaciones independientes por cada categoría. Esto permitiría realizar análisis estadísticos y evaluar la posibilidad de desarrollar un método automático de clasificación de señales.


% ============================================================
% REFERENCIAS
% ============================================================

\begin{thebibliography}{00}

\bibitem{oppenheim}
A. V. Oppenheim, A. S. Willsky y S. H. Nawab,
\textit{Signals and Systems},
2nd ed., Prentice Hall, 1996.

\bibitem{lyons}
R. G. Lyons,
\textit{Understanding Digital Signal Processing},
3rd ed., Pearson, 2011.

\bibitem{numpy}
C. R. Harris et al.,
``Array programming with NumPy,''
\textit{Nature}, vol. 585, pp. 357--362, 2020.

\bibitem{scipy}
P. Virtanen et al.,
``SciPy 1.0: Fundamental Algorithms for Scientific Computing in Python,''
\textit{Nature Methods}, vol. 17, pp. 261--272, 2020.

\end{thebibliography}